A seguir será tratado sobre o planejamento de trajetória autônomo, onde atualmente as espaçonaves utilizam para minimizar riscos de colisão, mesmos em cenários de falha parcial.

\subsection{Métodos Autônomos de Planejamento de Trajetória e Controle de Guiagem}

Os métodos autônomos de guiagem ou \textit{Autonomous Guidance Methods}  (AGMs), desde de 2016 têm sido estudados para melhorar a autonomia de voo em diferentes situações, esses métodos geram comandos de guiagem em tempo real e garantem que a trajetória atenda às restrições dinâmicas do veículo durante todas as fases do voo. Essa abordagem é diferente das abordagens tradicionais baseadas em trajetórias pré-planejadas, o diferencial é que os AGMs permitem maior flexibilidade ao lidar com restrições não lineares e variáveis durante o voo \citep{song2023}.

Inicialmente, dentre os métodos existentes, destacam-se abordagens como a \textit{Computational Guidance}, onde as leis de controle são substituídas por algoritmos baseados em modelos ou dados. Essa tipo de método elimina a necessidade de ajustes prévios e planejamento extensivo \textit{offline}. Outro modelo relevante é o \textit{Model-Based Real-Time Optimization}, que otimiza a trajetória de forma contínua, lidando com restrições complexas \citep{song2023}. Por fim, a \textit{Autonomous Mission Planning} detém o maior nível de autonomia e permite que o próprio sistema gerencie a missão em tempo real.

Conforme explica \citep{song2023}, outros conceitos que devem ser mencionados são a \textit{Adaptive Guidance}, que ajusta os parâmetros do sistema de acordo com as condições de voo e aprendizado contínuo, e a \textit{Adaptive Optimal Guidance}, que utiliza otimização numérica para compensar falhas de desempenho em estágios anteriores do lançamento. O conceito de \textit{Fault-Tolerant Guidance} também tem sido estudado, sendo ideal para cenários em que a trajetória nominal não pode ser alcançada devido a falhas no sistema de propulsão \citep{song2023}.

A aplicação desses métodos para aplicações espaciais exige considerações diferenciadas quando comparados a sistemas autônomos de navegação terrestre ou aeronáutica, pois no caso de veículos de lançamento, a não-linearidade da dinâmica de voo, as restrições gravitacionais e a variação significativa da massa ao longo da missão impõem desafios adicionais para a implementação de guiagem autônoma. Os AGMs têm demonstrado capacidade de solucionar esses desafios, trazendo segurança, robustez da navegação e planejamentos eficientes ao longo do voo \citep{song2023}.
